%-------------------------------------------------------------------------------
% Lettre d'accompagnement une page (français)
% Personnalisation : \RecipientName, \CompanyName ; \LetterSubject, \RoleTitle ;
%   \WhyCompany (1--2 phrases ; vide = paragraphe omis).
% Compilation : Set-Location letters/fr; latexmk -xelatex -interaction=nonstopmode cover-letter.tex
%-------------------------------------------------------------------------------
\documentclass[11pt, a4paper]{../../cv-fr/russell}

\geometry{left=1.4cm, top=1.0cm, right=1.4cm, bottom=2.2cm, footskip=.5cm}

\usepackage[normalem]{ulem}
\emergencystretch=3em
\hbadness=10000
\hfuzz=0.5pt

\fontdir[../../cv-fr/fonts/]

\colorlet{russell}{russell-black}
\setbool{acvSectionColorHighlight}{true}
\renewcommand{\acvHeaderSocialSep}{\quad\textbar\quad}

\newcommand{\RecipientName}{Responsable du recrutement}
\newcommand{\CompanyName}{}
\newcommand{\LetterSubject}{Objet : Candidature --- Ingénieur logiciel}
\newcommand{\RoleTitle}{des postes en plateforme, en infographie ou en développement logiciel généraliste (y compris sécurité et cloud)}
\newcommand{\WhyCompany}{}

\name{Antoine}{Boucher}
\position{Ingénieur logiciel --- Plateforme et infographie}
\mobile{+14384918831}
\email{antoine@antoineboucher.info}
\homepage{antoineboucher.info}
\github{antoinebou12}
\linkedin{antoineboucher12}

\begin{document}

\makecvheader[C]
\makecvfooter
  {\today}
  {}
  {\thepage}

\recipient{\RecipientName}{\ifdefempty{\CompanyName}{\vphantom{A}}{\CompanyName}}
\letterdate{\today}
\lettertitle{\LetterSubject}
\letteropening{Madame, Monsieur,}

\makelettertitle

\begin{cvletter}
\begin{spacing}{1.10}

Je termine une maîtrise en génie des technologies de l'information à l'École de technologie supérieure. Mon parcours couvre assez largement le développement logiciel : services dorsaux, interfaces web, automatisation, sécurité applicative, déploiement infonuagique, outils internes et programmation C++. Je me décris comme généraliste : j'apprends vite et je m'adapte au besoin de l'équipe. Je vise \RoleTitle. Mon curriculum vitæ détaille le parcours ; d'autres précisions figurent sur \href{https://antoineboucher.info}{antoineboucher.info}.

À l'IMC2, j'ai été auxiliaire de recherche à temps partiel (mars 2025 à février 2026) sur une plateforme de cybersécurité sur GCP. J'ai participé au développement et au renforcement de services Java, à l'automatisation de tests, au déploiement d'infrastructure avec Terraform et à l'exploitation de charges conteneurisées, ainsi qu'à l'intégration des compilations Android dans le même chemin de livraison que les services dorsaux. J'ai aussi contribué à des outils pour analystes : notation, visualisation et priorisation de risques cyber à grande échelle. J'y ai touché à plusieurs facettes d'un même produit --- backend, sécurité, données, visualisation, livraison et exploitation --- sans me limiter à un seul rôle technique.

J'ai été auxiliaire d'enseignement DevSecOps à Polytechnique Montréal (cours LOG8100), et auxiliaire pour plusieurs cours à l'ÉTS. J'ai accompagné des équipes sur GitLab, Docker, Kubernetes, Azure, les bases de données, le développement web et les pratiques sécuritaires. J'aime poser un problème clairement, débloquer les autres et faire avancer un livrable sous contrainte de temps.

Auparavant : infrastructure infonuagique chez IONODES, services Kotlin et Spring chez Intact, mandats web avec Django, Vue, API et observabilité. Ces expériences m'ont donné une base solide en développement backend et frontend, intégration, déploiement et maintenance.

Dans le cadre de la maîtrise, j'approfondis C++, OpenGL et Blender sur un projet de simulation d'usure de surface en temps réel dans l'espace texture, mêlant simulation, GPU et rendu interactif. Ce travail correspond à mon intérêt pour les systèmes exigeants en performance, la fiabilité et la qualité logicielle.

\ifdefempty{\WhyCompany}{}{%
\WhyCompany\par
}

Je publie aussi quelques projets open source, dont \textbf{\href{https://github.com/antoinebou12/uml-mcp}{uml-mcp}}, un outil lié à la génération de diagrammes UML. Je serais heureux d'échanger sur la manière dont mon parcours pourrait soutenir vos projets, votre plateforme ou vos opérations logicielles.

\end{spacing}
\end{cvletter}

\letterclosing{Cordialement,}
\letterenclosure[Pièce jointe]{Curriculum vitæ --- \href{https://antoineboucher.info}{antoineboucher.info} --- Download CV / Télécharger le CV}

\makeletterclosing

\end{document}
